\documentclass[11pt,a4paper]{article}
\usepackage{patient_robot_instructions}

\title{%
  \vspace{-1.5cm}
  {\Large\bfseries Patient-Robot Instructions: Physical AI Oncology Trials}\\[0.3em]
  {\normalsize Instructional Illustrations and Procedural Guidance for 10 Robot Types}
}
\author{%
  Kevin Kawchak\thanks{CEO, ChemicalQDevice. ORCID: \url{https://orcid.org/0009-0007-5457-8667}. Email: kevink@chemicalqdevice.com}\\
  \small DOI: \href{https://doi.org/10.5281/zenodo.18810541}{\texttt{10.5281/zenodo.18810541}}\\[0.2em]
  \small February 28, 2026
}
\date{}

\begin{document}

% ══════════════════════════════════════════════════════════════
% Each page is a self-contained instructional sheet for one robot type.
% Pages 5 and 6 feature pediatric patients.
% Robot types are ordered by frequency of use in physical AI oncology trials.
% Illustrations are black-and-white TikZ portraits rendered with Cairo.
%
% ISO standards considered:
%   - ISO 15223-1, ISO 20417 (medical device symbols and information)
%   - ISO 7000, IEC 60417 (graphical symbols for equipment)
%   - ISO 7010, ISO 3864-1 (safety signs and pictograms)
% ══════════════════════════════════════════════════════════════

% ──────────────────────────────────────────────────────────────
% PAGE 1: Surgical Robots
% ──────────────────────────────────────────────────────────────
\newpage
\thispagestyle{fancy}
\section*{Patient-Robot Instructions: Physical AI Oncology Trials -- Surgical Robots}

\begin{center}
\textit{Adult Oncology Trial Setting}

\vspace{0.5em}
% [Illustration: Adult patient lying on operating table. Multi-arm surgical robot
%  (da Vinci style) positioned overhead with 3--4 articulating arms. Camera arm
%  centered, instrument arms flanking. Patient in hospital gown, arms on rests.]
\fbox{\parbox{0.85\textwidth}{\centering\vspace{4cm}
\textbf{[TikZ/Cairo Illustration: Surgical Robot with Patient]}\\
Adult patient on operating table; multi-arm surgical robot overhead\\
ISO 7010 Caution Symbol \isowarning
\vspace{4cm}}}
\end{center}

\subsection*{1. Preparation at Home}
\begin{enumerate}
  \item Fast for 8 hours before your scheduled procedure (no food or drink).
  \item Shower with antibacterial soap the evening before and morning of.
  \item Wear loose, comfortable clothing; leave jewelry and valuables at home.
  \item Bring your photo ID, insurance card, and signed consent forms.
  \item Take prescribed pre-operative medications as directed by your oncologist.
\end{enumerate}

\subsection*{2. Entering the Room}
\begin{enumerate}
  \item A staff member will guide you to the pre-op area; the robot will already be in the OR.
  \item Change into the hospital gown provided and place belongings in the labeled bag.
  \item Lie on the operating table face-up when instructed.
  \item Place both arms on the padded arm rests, palms facing up.
  \item Remain still while the anesthesiologist administers sedation (approx.\ 2~min).
\end{enumerate}

\subsection*{3. During the Interaction}
\begin{enumerate}
  \item The surgical robot has 3--4 arms that will operate through small incisions (8--12~mm).
  \item You will be under general anesthesia and will not feel or need to do anything.
  \item Estimated procedure time: 45--180 minutes depending on tumor location.
  \item The surgeon controls the robot from a console nearby; the robot does not act autonomously.
  \item Robotic instruments include a camera, graspers, and cautery tools.
\end{enumerate}

\subsection*{4. Concluding the Session}
\begin{enumerate}
  \item After the procedure, the robot arms are retracted and incisions are closed.
  \item You will be moved to the recovery room; expect to wake in 15--30 minutes.
  \item A nurse will check your vitals every 15 minutes for the first hour.
  \item You may feel mild soreness at incision sites; request pain management as needed.
\end{enumerate}

\subsection*{5. At Home \& Follow-Up}
\begin{enumerate}
  \item Keep incision sites clean and dry for 48 hours; change dressings as instructed.
  \item Attend your post-operative follow-up appointment in 7--10 days.
  \item Report any fever above 38.5\,$^\circ$C, excessive bleeding, or unusual swelling immediately.
  \item Resume light activity after 1--2 weeks; avoid heavy lifting for 4--6 weeks.
  \item Your next robotic-assisted session (if needed) will be scheduled at follow-up.
\end{enumerate}

\vfill
\noindent\rule{\textwidth}{0.4pt}\\
{\scriptsize Sources: intuitivesurgical.com/davinci \cite{intuitive2026davinci}; \cite{sheetz2020trends}; \cite{moglia2021systematic}}

% ──────────────────────────────────────────────────────────────
% PAGE 2: Cobots
% ──────────────────────────────────────────────────────────────
\newpage
\thispagestyle{fancy}
\section*{Patient-Robot Instructions: Physical AI Oncology Trials -- Cobots}

\begin{center}
\textit{Adult Oncology Trial Setting}

\vspace{0.5em}
\fbox{\parbox{0.85\textwidth}{\centering\vspace{4cm}
\textbf{[TikZ/Cairo Illustration: Cobot with Patient]}\\
Adult patient seated in exam chair; 7-DOF collaborative robot arm on mobile cart\\
ISO 15223-1 Keep Still Symbol \isokeepstill
\vspace{4cm}}}
\end{center}

\subsection*{1. Preparation at Home}
\begin{enumerate}
  \item Eat a light meal 2 hours before your appointment to maintain comfort.
  \item Wear a short-sleeved shirt or top that allows arm and shoulder access.
  \item Review your medication list; bring any prescribed topical numbing cream.
  \item Arrive 15 minutes early to complete check-in paperwork.
  \item Remove watches, bracelets, and rings from the arm the cobot will assist.
\end{enumerate}

\subsection*{2. Entering the Room}
\begin{enumerate}
  \item Enter the procedure room; the cobot is mounted on a mobile cart beside the exam chair.
  \item Sit in the adjustable exam chair and place your arm on the padded arm rest.
  \item The cobot will be in standby mode with its arm retracted and a green status light on.
  \item A technician will verbally confirm the procedure and activate the cobot.
  \item Place your hand palm-down on the marked target area on the arm rest.
\end{enumerate}

\subsection*{3. During the Interaction}
\begin{enumerate}
  \item The cobot arm moves slowly and gently (max speed 250~mm/s); you may feel light pressure.
  \item If you feel discomfort, say ``Stop'' clearly; the cobot will halt within 0.5~seconds.
  \item Estimated interaction time: 10--25 minutes for sample collection or positioning tasks.
  \item Keep your arm relaxed and still; the cobot uses force sensors to avoid excess pressure.
  \item The cobot may reposition your arm 2--3 times; follow verbal audio prompts.
\end{enumerate}

\subsection*{4. Concluding the Session}
\begin{enumerate}
  \item The cobot arm will retract to its home position when the task is complete.
  \item Remain seated for 2 minutes while the technician reviews results on screen.
  \item A bandage will be applied to any sample collection site.
  \item You will receive a printed summary of the session before leaving.
\end{enumerate}

\subsection*{5. At Home \& Follow-Up}
\begin{enumerate}
  \item No special home care is usually required after cobot-assisted procedures.
  \item If a sample was collected, results will be available in 3--5 business days.
  \item Monitor the site for 24 hours; apply gentle pressure if any minor bleeding.
  \item Your next cobot session will be scheduled during your regular oncology visit.
\end{enumerate}

\vfill
\noindent\rule{\textwidth}{0.4pt}\\
{\scriptsize Sources: franka.de/research \cite{franka2026research}; \cite{lamon2023collaborative}; \cite{haddadin2022collaborative}}

% ──────────────────────────────────────────────────────────────
% PAGE 3: Radiotherapy Patient-Positioning Robots
% ──────────────────────────────────────────────────────────────
\newpage
\thispagestyle{fancy}
\section*{Patient-Robot Instructions: Physical AI Oncology Trials -- Radiotherapy Patient-Positioning Robots}

\begin{center}
\textit{Adult Oncology Trial Setting}

\vspace{0.5em}
\fbox{\parbox{0.85\textwidth}{\centering\vspace{4cm}
\textbf{[TikZ/Cairo Illustration: Radiotherapy Positioning Robot with Patient]}\\
Adult patient lying on robotic 6-DOF couch; gantry ring above; laser crosshairs\\
ISO 7010 W003 Ionizing Radiation Symbol \isoradiation
\vspace{4cm}}}
\end{center}

\subsection*{1. Preparation at Home}
\begin{enumerate}
  \item Wear comfortable clothing without metal (zippers, snaps, underwire).
  \item Do not apply lotion, deodorant, or powder to the treatment area.
  \item Eat and drink normally unless otherwise directed by your radiation oncologist.
  \item Take anxiety-reducing medication 30 min before if prescribed.
  \item Use the restroom before your session; a full bladder may be required for some sites.
\end{enumerate}

\subsection*{2. Entering the Room}
\begin{enumerate}
  \item Enter the treatment room; the robotic couch is the flat table in the center.
  \item Lie on the robotic couch in the position demonstrated during your simulation session.
  \item Place your arms in the arm cradle or above your head as instructed.
  \item A custom immobilization mask or mold may be placed to hold you steady.
  \item The couch will adjust automatically using 6 degrees of freedom (6-DOF).
\end{enumerate}

\subsection*{3. During the Interaction}
\begin{enumerate}
  \item The robotic couch shifts in small increments (1--3~mm) to align you with the beam.
  \item You will hear mechanical sounds as the couch moves; this is normal.
  \item Breathe normally and remain completely still during alignment (approx.\ 5--8~min).
  \item Laser crosshairs will project onto your skin to verify position; do not block them.
  \item Estimated total session time: 15--30 minutes including setup and treatment.
\end{enumerate}

\subsection*{4. Concluding the Session}
\begin{enumerate}
  \item After the beam delivery ends, the couch will lower to exit height.
  \item Wait for the therapist to remove any immobilization devices before sitting up.
  \item Sit up slowly to avoid dizziness and swing your legs off the side.
  \item The therapist will confirm your next session date and time.
\end{enumerate}

\subsection*{5. At Home \& Follow-Up}
\begin{enumerate}
  \item Gently wash the treatment area with mild soap; do not scrub.
  \item Apply recommended moisturizer to the treatment area if skin redness develops.
  \item Attend all scheduled daily sessions (typically Mon--Fri for 3--7 weeks).
  \item Report skin changes, pain, or fatigue to your radiation oncology team.
\end{enumerate}

\vfill
\noindent\rule{\textwidth}{0.4pt}\\
{\scriptsize Sources: accuray.com/cyberknife \cite{accuray2026cyberknife}; \cite{hoisak2022robotic}; \cite{verellen2003robotic}}

% ──────────────────────────────────────────────────────────────
% PAGE 4: Robotic Needle-Placement Systems
% ──────────────────────────────────────────────────────────────
\newpage
\thispagestyle{fancy}
\section*{Patient-Robot Instructions: Physical AI Oncology Trials -- Robotic Needle-Placement Systems}

\begin{center}
\textit{Adult Oncology Trial Setting}

\vspace{0.5em}
\fbox{\parbox{0.85\textwidth}{\centering\vspace{4cm}
\textbf{[TikZ/Cairo Illustration: Robotic Needle-Placement System with Patient]}\\
Adult patient on imaging table; articulating needle-guide robot with display\\
ISO 7010 Caution Symbol \isowarning
\vspace{4cm}}}
\end{center}

\subsection*{1. Preparation at Home}
\begin{enumerate}
  \item Fast for 4 hours before the procedure (clear liquids may be allowed).
  \item Inform the team of any blood-thinning medications; some may need to be paused.
  \item Wear loose clothing that provides access to the biopsy or treatment area.
  \item Arrange for someone to drive you home if sedation will be used.
  \item Bring prior imaging (CT/MRI scans) on CD or USB if requested.
\end{enumerate}

\subsection*{2. Entering the Room}
\begin{enumerate}
  \item Enter the procedure room; the robotic needle system is beside the imaging table.
  \item Lie on the imaging table in the position indicated (face-up or face-down).
  \item Place both hands at your sides or above your head as directed.
  \item A local anesthetic will be applied to the needle insertion site (1--2~min to numb).
  \item The robot's articulating arm will position the needle guide over the target.
\end{enumerate}

\subsection*{3. During the Interaction}
\begin{enumerate}
  \item The robot aligns the needle to sub-millimeter accuracy using real-time imaging.
  \item You may feel pressure but should not feel sharp pain; alert staff if you do.
  \item Hold your breath briefly (5--10~seconds) when instructed for image capture.
  \item The needle insertion takes 1--3~minutes; the robot maintains steady guidance.
  \item Estimated total procedure time: 20--45~minutes including imaging and sampling.
\end{enumerate}

\subsection*{4. Concluding the Session}
\begin{enumerate}
  \item The needle is withdrawn and gentle pressure is applied for 3--5~minutes.
  \item A sterile bandage is placed over the insertion site.
  \item Remain lying down for 15--30~minutes for post-procedure monitoring.
  \item A nurse will check your blood pressure and verify no immediate complications.
\end{enumerate}

\subsection*{5. At Home \& Follow-Up}
\begin{enumerate}
  \item Keep the bandage dry for 24 hours; replace with a clean one after.
  \item Avoid strenuous activity or heavy lifting for 48 hours.
  \item Biopsy results are typically available within 5--7 business days.
  \item Contact your care team if you experience increasing pain, swelling, or bleeding.
  \item A follow-up appointment will be scheduled to review results and next steps.
\end{enumerate}

\vfill
\noindent\rule{\textwidth}{0.4pt}\\
{\scriptsize Sources: \cite{walsh2022robotic}; \cite{li2023needle}; ISO 15223-1 \cite{iso15223}}

% ──────────────────────────────────────────────────────────────
% PAGE 5: Social Companion Robots (PEDIATRIC)
% ──────────────────────────────────────────────────────────────
\newpage
\thispagestyle{fancy}
\section*{Patient-Robot Instructions: Physical AI Oncology Trials -- Social Companion Robots}

\begin{center}
\textit{Pediatric Oncology Trial Setting}

\vspace{0.5em}
\fbox{\parbox{0.85\textwidth}{\centering\vspace{4cm}
\textbf{[TikZ/Cairo Illustration: Social Companion Robot with Child Patient]}\\
Child patient interacting with friendly Pepper-style robot; chest screen with heart icon\\
ISO 15223-1 Keep Still Symbol \isokeepstill
\vspace{4cm}}}
\end{center}

\subsection*{1. Preparation at Home}
\begin{enumerate}
  \item Parents/guardians: explain to your child that they will meet a friendly robot.
  \item Bring your child's favorite small comfort item (stuffed animal, blanket).
  \item Ensure your child has eaten a light snack and is well-hydrated.
  \item Dress your child in comfortable, easy-to-move-in clothing.
  \item Review the picture book provided at your last visit showing the robot.
\end{enumerate}

\subsection*{2. Entering the Room}
\begin{enumerate}
  \item Your child enters the room alone (you may watch through the window).
  \item The companion robot (about 1.2~m tall) will greet your child by name.
  \item The robot's screen displays a friendly face and a welcome message.
  \item Encourage your child to wave or say ``Hello'' to the robot.
  \item The robot will extend its hand for a gentle high-five if your child is comfortable.
\end{enumerate}

\subsection*{3. During the Interaction}
\begin{enumerate}
  \item The robot guides your child through breathing exercises with visual cues on its screen.
  \item Your child can talk to the robot; it responds with simple, comforting phrases.
  \item Interactive games on the robot's chest screen help distract during waiting (5--10~min).
  \item The robot demonstrates how to sit still using fun animations (e.g., ``freeze like a statue'').
  \item Estimated interaction time: 10--20 minutes before and/or during the clinical procedure.
\end{enumerate}

\subsection*{4. Concluding the Session}
\begin{enumerate}
  \item The robot will say goodbye and offer a virtual sticker on its screen.
  \item Your child can give the robot a high-five or wave goodbye.
  \item The robot plays a short congratulatory song (15~seconds).
  \item A staff member will escort your child back to you.
\end{enumerate}

\subsection*{5. At Home \& Follow-Up}
\begin{enumerate}
  \item Ask your child about their experience to reinforce positive associations.
  \item The companion robot session notes are shared with your child's oncology team.
  \item Future visits will include the same robot to build familiarity and trust.
  \item Report any anxiety or distress your child expresses about the robot to the care team.
\end{enumerate}

\vfill
\noindent\rule{\textwidth}{0.4pt}\\
{\scriptsize Sources: softbankrobotics.com/pepper \cite{softbank2026pepper}; \cite{logan2019social}; \cite{rhee2023companion}}

% ──────────────────────────────────────────────────────────────
% PAGE 6: Humanoids (PEDIATRIC)
% ──────────────────────────────────────────────────────────────
\newpage
\thispagestyle{fancy}
\section*{Patient-Robot Instructions: Physical AI Oncology Trials -- Humanoids}

\begin{center}
\textit{Pediatric Oncology Trial Setting}

\vspace{0.5em}
\fbox{\parbox{0.85\textwidth}{\centering\vspace{4cm}
\textbf{[TikZ/Cairo Illustration: Humanoid Robot with Child Patient]}\\
Child patient walking alongside bipedal humanoid robot; robot at child's eye level\\
ISO 15223-1 Keep Still Symbol \isokeepstill
\vspace{4cm}}}
\end{center}

\subsection*{1. Preparation at Home}
\begin{enumerate}
  \item Parents/guardians: show your child the humanoid robot introduction video sent by email.
  \item Pack a small snack and water bottle for after the session.
  \item Dress your child in comfortable clothes and non-slip shoes for walking activities.
  \item Arrive 10 minutes early so your child can see the robot through the window first.
  \item Inform the team of your child's preferred name and any communication needs.
\end{enumerate}

\subsection*{2. Entering the Room}
\begin{enumerate}
  \item Your child enters the activity room; the humanoid robot (about 1.5~m tall) stands inside.
  \item The robot waves and introduces itself by name in a child-friendly voice.
  \item The robot takes one slow step forward and stops, allowing your child to approach.
  \item Your child may hold the robot's hand if comfortable; the grip is very gentle.
  \item The robot kneels down to your child's eye level for face-to-face interaction.
\end{enumerate}

\subsection*{3. During the Interaction}
\begin{enumerate}
  \item The robot walks alongside your child at their pace (0.5--1.0~m/s maximum).
  \item Together they practice the route to the treatment room (approx.\ 3--5~minutes).
  \item The robot demonstrates deep breathing: ``Breathe in 1-2-3, out 1-2-3.'' Your child copies.
  \item Your child can ask the robot questions; it answers in simple, reassuring language.
  \item Estimated session time: 15--25 minutes including walking and guided preparation.
\end{enumerate}

\subsection*{4. Concluding the Session}
\begin{enumerate}
  \item The robot walks your child back to the waiting area.
  \item It gives a gentle wave and says, ``Great job! See you next time.''
  \item The robot returns to its standby position in the activity room.
  \item A care team member brings your child to you with a session summary card.
\end{enumerate}

\subsection*{5. At Home \& Follow-Up}
\begin{enumerate}
  \item Discuss with your child what they liked about the humanoid robot visit.
  \item The session summary is added to your child's electronic health record.
  \item The same humanoid robot will be used for consistency at future visits.
  \item Contact the pediatric oncology team if your child has questions or concerns.
\end{enumerate}

\vfill
\noindent\rule{\textwidth}{0.4pt}\\
{\scriptsize Sources: bostondynamics.com/atlas \cite{bostondynamics2026atlas}; \cite{darvish2023humanoid}}

% ──────────────────────────────────────────────────────────────
% PAGE 7: Radiotherapy Motion-Management / Tracking Robots
% ──────────────────────────────────────────────────────────────
\newpage
\thispagestyle{fancy}
\section*{Patient-Robot Instructions: Physical AI Oncology Trials -- Radiotherapy Motion-Management / Tracking Robots}

\begin{center}
\textit{Adult Oncology Trial Setting}

\vspace{0.5em}
\fbox{\parbox{0.85\textwidth}{\centering\vspace{4cm}
\textbf{[TikZ/Cairo Illustration: Motion-Tracking Robot with Patient]}\\
Adult patient on couch; stereo camera array overhead; reflective markers on chest\\
ISO 7010 W003 Ionizing Radiation Symbol \isoradiation
\vspace{4cm}}}
\end{center}

\subsection*{1. Preparation at Home}
\begin{enumerate}
  \item Wear comfortable clothing without metal; a hospital gown may be provided.
  \item Do not eat a large meal within 1 hour of your session to reduce diaphragm movement.
  \item Practice the breathing pattern taught during your simulation: inhale 3~sec, exhale 3~sec.
  \item Avoid caffeine on the day of treatment, as it may increase involuntary movement.
  \item Bring your treatment schedule card for session verification.
\end{enumerate}

\subsection*{2. Entering the Room}
\begin{enumerate}
  \item Enter the treatment vault; the motion-tracking system has cameras mounted overhead.
  \item Lie on the treatment couch and position your arms as practiced in simulation.
  \item A technician places a small reflective marker block on your chest or abdomen.
  \item Breathe normally; the tracking cameras detect the markers and display your breathing trace.
  \item The system calibrates to your breathing pattern in approximately 1--2~minutes.
\end{enumerate}

\subsection*{3. During the Interaction}
\begin{enumerate}
  \item Follow the breathing guide on the ceiling-mounted screen: inhale when the bar rises.
  \item The system tracks your tumor position in real time; the beam activates only when aligned.
  \item You will hear clicking sounds as the beam gates on and off; this is expected.
  \item Remain as still as possible between breaths; the tracking tolerance is 2--3~mm.
  \item Estimated treatment delivery time: 10--20~minutes once tracking is calibrated.
\end{enumerate}

\subsection*{4. Concluding the Session}
\begin{enumerate}
  \item After the last beam, the tracking system deactivates and cameras power down.
  \item The technician removes the marker block from your chest.
  \item Sit up slowly and wait for the therapist to confirm session completion.
  \item Your breathing trace data is saved for comparison at future sessions.
\end{enumerate}

\subsection*{5. At Home \& Follow-Up}
\begin{enumerate}
  \item Practice your breathing exercises at home for 5 minutes daily between sessions.
  \item Attend all scheduled treatment sessions to maintain the motion-management protocol.
  \item Report any new breathing difficulties or chest discomfort to your care team.
  \item Skin care for the treatment area follows the same guidelines as standard radiotherapy.
\end{enumerate}

\vfill
\noindent\rule{\textwidth}{0.4pt}\\
{\scriptsize Sources: varian.com/respiratory-gating \cite{varian2026gating}; \cite{bertholet2019motion}; \cite{keall2023tumor}}

% ──────────────────────────────────────────────────────────────
% PAGE 8: Imaging Assistant Robots
% ──────────────────────────────────────────────────────────────
\newpage
\thispagestyle{fancy}
\section*{Patient-Robot Instructions: Physical AI Oncology Trials -- Imaging Assistant Robots}

\begin{center}
\textit{Adult Oncology Trial Setting}

\vspace{0.5em}
\fbox{\parbox{0.85\textwidth}{\centering\vspace{4cm}
\textbf{[TikZ/Cairo Illustration: Imaging Assistant Robot with Patient]}\\
Adult patient on exam bed; robotic ultrasound arm with force-controlled probe holder\\
ISO 15223-1 Keep Still Symbol \isokeepstill
\vspace{4cm}}}
\end{center}

\subsection*{1. Preparation at Home}
\begin{enumerate}
  \item Drink 500~mL of water 1 hour before your appointment for ultrasound imaging clarity.
  \item Wear a two-piece outfit for easy access to the imaging area.
  \item Do not apply oils or lotion to the area being imaged.
  \item Bring previous imaging reports if this is a follow-up scan.
  \item Arrive 10 minutes early to complete a brief imaging questionnaire.
\end{enumerate}

\subsection*{2. Entering the Room}
\begin{enumerate}
  \item Enter the imaging room; the robotic ultrasound system is beside the exam bed.
  \item Lie on the exam bed and expose the area to be imaged.
  \item The robotic arm holds the ultrasound probe and will approach your body slowly.
  \item The technician applies ultrasound gel to the imaging area (it may feel cool).
  \item Place your hands behind your head or at your sides as instructed.
\end{enumerate}

\subsection*{3. During the Interaction}
\begin{enumerate}
  \item The robotic probe applies consistent, gentle pressure (1--3~N) on your skin.
  \item Breathe normally; the robot compensates for your breathing motion automatically.
  \item The robot scans methodically in a grid pattern for complete coverage of the target area.
  \item You may hear the motor of the robotic arm; this is normal operation.
  \item Estimated imaging time: 10--20 minutes depending on the area and scan complexity.
\end{enumerate}

\subsection*{4. Concluding the Session}
\begin{enumerate}
  \item The robotic arm retracts and the probe is lifted from your skin.
  \item The technician wipes excess gel from the imaged area with a warm cloth.
  \item Preliminary images appear on the display for immediate quality check.
  \item You may sit up and get dressed; there is no recovery period needed.
\end{enumerate}

\subsection*{5. At Home \& Follow-Up}
\begin{enumerate}
  \item Imaging results are reviewed by a radiologist within 1--3 business days.
  \item Your oncologist will discuss findings at your next scheduled appointment.
  \item No special home care is needed after robotic ultrasound imaging.
  \item Future robotic imaging sessions maintain the same scan protocol for comparison.
\end{enumerate}

\vfill
\noindent\rule{\textwidth}{0.4pt}\\
{\scriptsize Sources: \cite{von2022robotic}; \cite{jiang2023imaging}; ISO 20417 \cite{iso20417}}

% ──────────────────────────────────────────────────────────────
% PAGE 9: Steerable Needle / Needle-Steering Robots
% ──────────────────────────────────────────────────────────────
\newpage
\thispagestyle{fancy}
\section*{Patient-Robot Instructions: Physical AI Oncology Trials -- Steerable Needle / Needle-Steering Robots}

\begin{center}
\textit{Adult Oncology Trial Setting}

\vspace{0.5em}
\fbox{\parbox{0.85\textwidth}{\centering\vspace{4cm}
\textbf{[TikZ/Cairo Illustration: Steerable Needle Robot with Patient]}\\
Adult patient on table; needle-steering robot with curved needle path to target\\
ISO 7010 Caution Symbol \isowarning
\vspace{4cm}}}
\end{center}

\subsection*{1. Preparation at Home}
\begin{enumerate}
  \item Fast for 6 hours before the procedure if sedation is planned.
  \item Inform the team of all current medications, especially blood thinners.
  \item An MRI compatibility screening form must be completed if MRI-guided.
  \item Remove all metal objects, credit cards, and electronic devices.
  \item Arrange transportation home; you should not drive after sedation.
\end{enumerate}

\subsection*{2. Entering the Room}
\begin{enumerate}
  \item Enter the interventional suite; the needle-steering robot is mounted to the imaging table.
  \item Lie on the table in the position indicated; pillows support your comfort.
  \item The area around the insertion site is cleaned and draped with sterile towels.
  \item A local anesthetic is injected at the entry point (brief 5-second sting).
  \item The robot's needle guide is positioned at the skin entry point.
\end{enumerate}

\subsection*{3. During the Interaction}
\begin{enumerate}
  \item The steerable needle curves through tissue along a pre-planned path to the target.
  \item You may feel deep pressure but should not feel sharp pain; speak up if you do.
  \item Real-time MRI or CT images confirm the needle tip position every 5--10~seconds.
  \item Hold still; the robot compensates for small movements but large shifts require re-planning.
  \item Estimated procedure time: 30--60 minutes for needle placement and treatment delivery.
\end{enumerate}

\subsection*{4. Concluding the Session}
\begin{enumerate}
  \item The steerable needle is slowly withdrawn along its insertion path.
  \item Pressure is held at the entry site for 5--10 minutes.
  \item A sterile dressing is applied and secured with medical tape.
  \item You are monitored in recovery for 30--60 minutes before discharge.
\end{enumerate}

\subsection*{5. At Home \& Follow-Up}
\begin{enumerate}
  \item Rest at home for 24 hours; avoid bending or straining.
  \item Keep the dressing clean and dry for 48 hours.
  \item Take prescribed pain medication as directed (typically acetaminophen).
  \item Follow-up imaging is usually scheduled 1--2 weeks after the procedure.
  \item Call your care team for fever, worsening pain, or signs of infection.
\end{enumerate}

\vfill
\noindent\rule{\textwidth}{0.4pt}\\
{\scriptsize Sources: \cite{cowan2023steerable}; \cite{reed2023needle}; ISO 7010 \cite{iso7010}}

% ──────────────────────────────────────────────────────────────
% PAGE 10: Rehabilitation Exoskeletons / Robotic Gait Trainers
% ──────────────────────────────────────────────────────────────
\newpage
\thispagestyle{fancy}
\section*{Patient-Robot Instructions: Physical AI Oncology Trials -- Rehabilitation Exoskeletons / Robotic Gait Trainers}

\begin{center}
\textit{Adult Oncology Trial Setting}

\vspace{0.5em}
\fbox{\parbox{0.85\textwidth}{\centering\vspace{4cm}
\textbf{[TikZ/Cairo Illustration: Rehabilitation Exoskeleton with Patient]}\\
Adult patient wearing bilateral leg exoskeleton; gripping parallel bars\\
ISO 15223-1 Keep Still Symbol \isokeepstill
\vspace{4cm}}}
\end{center}

\subsection*{1. Preparation at Home}
\begin{enumerate}
  \item Wear fitted athletic clothing (leggings or joggers, T-shirt) and flat athletic shoes.
  \item Eat a moderate meal 1--2 hours before your session for sustained energy.
  \item Complete the pre-session mobility questionnaire sent to your phone.
  \item Bring a water bottle; hydration is important during physical rehabilitation.
  \item Inform the therapist of any new pain, swelling, or changes since your last session.
\end{enumerate}

\subsection*{2. Entering the Room}
\begin{enumerate}
  \item Enter the rehabilitation gym; the exoskeleton is hanging on a support frame between parallel bars.
  \item Stand with your back to the exoskeleton; a therapist secures the waist belt first.
  \item Thigh and shank straps are fastened snugly (you should feel firm but not tight pressure).
  \item Foot plates are adjusted to match your shoe size and secured with Velcro.
  \item Grip the parallel bars with both hands before the exoskeleton powers on.
\end{enumerate}

\subsection*{3. During the Interaction}
\begin{enumerate}
  \item The exoskeleton initiates a gentle standing motion; let it support your weight.
  \item Follow the audible ``step'' cues: the exoskeleton moves one leg, then the other.
  \item Walk between the parallel bars at the robot's guided pace (0.2--0.5~m/s).
  \item If fatigued, say ``Pause''; the exoskeleton stops and supports you in a standing position.
  \item Estimated active walking time: 15--30 minutes with rest breaks as needed.
\end{enumerate}

\subsection*{4. Concluding the Session}
\begin{enumerate}
  \item The exoskeleton slowly returns you to a standing rest position.
  \item The therapist unbuckles straps in reverse order: feet, shank, thigh, waist.
  \item Sit in the provided chair for 5 minutes; drink water and rest.
  \item The therapist records your session metrics (steps, distance, gait symmetry).
\end{enumerate}

\subsection*{5. At Home \& Follow-Up}
\begin{enumerate}
  \item Perform the prescribed home stretching exercises daily between sessions.
  \item Sessions are typically 2--3 times per week for 4--8 weeks.
  \item Track your fatigue, pain, and mobility in the patient diary provided.
  \item Progress metrics are reviewed at your monthly oncology rehabilitation appointment.
\end{enumerate}

\vfill
\noindent\rule{\textwidth}{0.4pt}\\
{\scriptsize Sources: \cite{molteni2023exoskeleton}; eksobionics.com/ekso-indego \cite{ekso2026indego}; ISO 15223-1 \cite{iso15223}}

\end{document}
